\chapter{Aljabar dan Persamaan}\label{ch:g9:algebra}

Huruf mewakili bilangan, dan berhitung dengan huruf membuktikan
kenyataan tentang \emph{semua} bilangan sekaligus. Bab ini melatih
\omterm{def:g9:algebra:expand}{penjabaran} dan pemfaktoran --- termasuk ketiga kesamaan yang akan
mengikutimu di seluruh matematika --- lalu menyelesaikan \omterm{def:g8:equations:def}{persamaan}
berderajat satu, \omterm{def:g8:equations:def}{persamaan} \omterm{def:g2:mult:def}{hasil kali}, dan pertidaksamaan sederhana,
selalu satu langkah beralasan setiap kalinya.

\section{Menjabarkan}

\begin{definition}[Menjabarkan]\label{def:g9:algebra:expand}
\emph{Menjabarkan}\index{penjabaran} berarti mengubah sebuah \omterm{def:g2:mult:def}{hasil
kali} menjadi sebuah \omterm{def:g1:addition:def}{jumlah}, dengan memakai sifat distributif:
\[
k(a + b) = ka + kb,
\qquad
(a + b)(c + d) = ac + ad + bc + bd .
\]
\end{definition}

\begin{example}\label{ex:g9:algebra:expand}
Tiap suku tanda kurung pertamanya mengalikan tiap suku yang kedua:
\begin{align*}
(x + 3)(2x + 5)
&= x \times 2x + x \times 5 + 3 \times 2x + 3 \times 5 \\
&= 2x^2 + 5x + 6x + 15 \\
&= 2x^2 + 11x + 15 .
\end{align*}
Tandanya ikut bepergian bersama sukunya:
\begin{align*}
(3x - 2)(x - 4)
&= 3x^2 - 12x - 2x + 8
= 3x^2 - 14x + 8 .
\end{align*}
\end{example}

\begin{theorem}[Ketiga kesamaannya]\label{thm:g9:algebra:identities}
Untuk semua bilangan $a$ dan $b$:
\[
(a+b)^2 = a^2 + 2ab + b^2,
\qquad
(a-b)^2 = a^2 - 2ab + b^2,
\qquad
(a+b)(a-b) = a^2 - b^2 .
\]
\end{theorem}

\begin{proof}
Jabarkan tiap ruas kirinya. Untuk yang pertama:
$(a+b)(a+b) = a^2 + ab + ba + b^2 = a^2 + 2ab + b^2$. Untuk yang
kedua, gantilah $b$ dengan $-b$ pada yang pertama. Untuk yang ketiga:
$(a+b)(a-b) = a^2 - ab + ba - b^2 = a^2 - b^2$ (suku silangnya saling
meniadakan).
\end{proof}

\begin{omfigure}
\begin{tikzpicture}[scale=0.60]
  \draw[thick] (0,0) rectangle (5,5);
  \draw[thick] (3.3,0) -- (3.3,5);
  \draw[thick] (0,3.3) -- (5,3.3);
  \fill[omDef!14] (0,0) rectangle (3.3,3.3);
  \fill[omThm!14] (3.3,3.3) rectangle (5,5);
  \fill[omProp!14] (3.3,0) rectangle (5,3.3);
  \fill[omProp!14] (0,3.3) rectangle (3.3,5);
  \node[font=\small] at (1.65,1.65) {$a^2$};
  \node[font=\small] at (4.15,4.15) {$b^2$};
  \node[font=\small] at (4.15,1.65) {$ab$};
  \node[font=\small] at (1.65,4.15) {$ab$};
  \node[below, font=\small] at (1.65,0) {$a$};
  \node[below, font=\small] at (4.15,0) {$b$};
  \node[left, font=\small] at (0,1.65) {$a$};
  \node[left, font=\small] at (0,4.15) {$b$};
\end{tikzpicture}

{\small Kesamaan $(a+b)^2 = a^2 + 2ab + b^2$, yang dilihat sebagai
luas: persegi besar bersisi $a+b$ tersusun dari dua persegi dan dua
persegi panjang yang sama. Melupakan ``$2ab$'' berarti melupakan kedua
persegi panjangnya!}
\end{omfigure}

\begin{example}\label{ex:g9:algebra:identities}
\[
(x + 5)^2 = x^2 + 10x + 25, \qquad
(3x - 1)^2 = 9x^2 - 6x + 1, \qquad
(2x + 7)(2x - 7) = 4x^2 - 49 .
\]
Berhitung di kepala: $102^2 = (100 + 2)^2 = 10000 + 400 + 4 = 10404$,
dan $98 \times 102 = (100-2)(100+2) = 10000 - 4 = 9996$.
\end{example}

\section{Memfaktorkan}

\begin{method}[Memfaktorkan]\label{met:g9:algebra:factor}
Untuk memfaktorkan sebuah bentuk, cobalah berurutan:
\begin{enumerate}
  \item \emph{faktor bersama}: $ka + kb = k(a + b)$, dengan $k$ boleh
        berupa tanda kurung juga;
  \item \emph{\omterm{ex:g1:subtraction:difference}{selisih} kuadrat}: $a^2 - b^2 = (a+b)(a-b)$;
  \item \emph{kuadrat sempurna}: $a^2 + 2ab + b^2 = (a+b)^2$ (dan
        dengan $-2ab$, menjadi $(a-b)^2$).
\end{enumerate}
Periksalah selalu dengan \omterm{def:g9:algebra:expand}{menjabarkan} hasilnya.
\end{method}

\begin{example}\label{ex:g9:algebra:factor}
\emph{Faktor bersama:}
$15x^2 + 10x = 5x \times 3x + 5x \times 2 = 5x(3x + 2)$.

\emph{Tanda kurung bersama:}
\[
(x+2)(x-1) + (x+2)(3x+4) = (x+2)\bigl[(x-1) + (3x+4)\bigr]
= (x+2)(4x+3).
\]

\emph{\omterm{ex:g1:subtraction:difference}{Selisih} kuadrat:}
$x^2 - 81 = (x+9)(x-9)$, dan
\[
(2x+1)^2 - 9 = (2x+1)^2 - 3^2 = (2x+1+3)(2x+1-3) = (2x+4)(2x-2).
\]
Tiap faktornya masih punya faktor bersama: $(2x+4)(2x-2) = 2(x+2)
\times 2(x-1) = 4(x+2)(x-1)$.
\end{example}

\section{Persamaan}

\begin{method}[Persamaan berderajat satu]\label{met:g9:algebra:firstdegree}
Untuk menyelesaikan \omterm{def:g8:equations:def}{persamaan} seperti $5x - 3 = 2x + 9$:
\begin{enumerate}
  \item kumpulkan suku $x$ pada satu ruas dan bilangannya pada ruas
        yang lain, dengan menambahkan besaran yang sama pada kedua
        ruasnya;
  \item sederhanakan tiap ruasnya;
  \item bagilah kedua ruasnya dengan koefisien $x$;
  \item periksalah penyelesaiannya pada \omterm{def:g8:equations:def}{persamaan} aslinya.
\end{enumerate}
\end{method}

\begin{example}\label{ex:g9:algebra:firstdegree}
\begin{align*}
5x - 3 &= 2x + 9 \\
5x - 2x &= 9 + 3 && \text{(tambahkan $3 - 2x$ pada kedua ruasnya)}\\
3x &= 12 \\
x &= 4 && \text{(bagilah dengan 3).}
\end{align*}
Periksa: $5 \times 4 - 3 = 17$ dan $2 \times 4 + 9 = 17$.
Penyelesaiannya adalah $4$.

Dengan tanda kurung, jabarkan lebih dahulu: $3(x - 2) = x + 8$ menjadi
$3x - 6 = x + 8$, jadi $2x = 14$ dan $x = 7$.
\end{example}

\begin{theorem}[Aturan hasil kali nol]\label{thm:g9:algebra:zeroproduct}
Sebuah \omterm{def:g2:mult:def}{hasil kali} bernilai \omterm{def:g1:counting:zero}{nol} jika dan hanya jika sedikitnya satu
faktornya \omterm{def:g1:counting:zero}{nol}:
\[
A \times B = 0
\quad\Longleftrightarrow\quad
A = 0 \ \text{ atau } \ B = 0 .
\]
\end{theorem}

\begin{proof}
Kalau sebuah faktornya \omterm{def:g1:counting:zero}{nol}, jelas \omterm{def:g2:mult:def}{hasil kalinya} \omterm{def:g1:counting:zero}{nol}. Sebaliknya kalau
$A B = 0$ dengan $A \neq 0$, membagi kedua ruasnya dengan $A$ memberi
$B = 0$.
\end{proof}

\begin{example}\label{ex:g9:algebra:zeroproduct}
Selesaikan $(x - 3)(2x + 10) = 0$: entah $x - 3 = 0$ entah
$2x + 10 = 0$, jadi penyelesaiannya adalah $3$ dan $-5$.

Selesaikan $x^2 - 49 = 0$: faktorkan lebih dahulu, $(x-7)(x+7) = 0$,
penyelesaiannya $7$ dan $-7$.

Selesaikan $x^2 = 3x$: jangan pernah membagi dengan $x$! Pindahkan
lalu faktorkan: $x^2 - 3x = x(x - 3) = 0$, penyelesaiannya $0$ dan $3$.
\end{example}

\section{Pertidaksamaan}

\begin{proposition}[Aturan untuk pertidaksamaan]\label{prop:g9:algebra:ineq}
Sebuah pertidaksamaan terpelihara ketika bilangan yang sama
ditambahkan pada kedua ruasnya, dan ketika kedua ruasnya dikalikan
dengan bilangan \emph{positif} yang sama. Mengalikan kedua ruasnya
dengan bilangan \emph{negatif} membalikkan tanda
pertidaksamaannya.
\end{proposition}

\begin{proof}
Kalau $a \leq b$ maka $b - a \geq 0$. Dengan menambahkan $c$:
$(b+c) - (a+c) = b - a \geq 0$, jadi $a + c \leq b + c$. Dengan
mengalikan dengan $c > 0$: $bc - ac = (b-a)c \geq 0$, jadi
$ac \leq bc$. Dengan mengalikan dengan $c < 0$:
$(b-a)c \leq 0$, jadi $ac \geq bc$: urutannya terbalik.
\end{proof}

\begin{example}\label{ex:g9:algebra:ineq}
Selesaikan $7 - 2x < 15$:
\begin{align*}
7 - 2x &< 15 \\
-2x &< 8 && \text{(kurangkan 7)} \\
x &> -4 && \text{(bagilah dengan $-2$: balikkan tandanya).}
\end{align*}
Penyelesaiannya adalah semua bilangan yang lebih besar daripada $-4$:
pada garis bilangan, sebuah titik berongga di $-4$ dan arsiran ke kanan.
\end{example}

\begin{omfigure}
\begin{tikzpicture}[scale=1.0]
  \draw[-Stealth] (-6.4,0) -- (2.6,0) node[below, font=\small] {$x$};
  \foreach \x in {-6,-5,-4,-3,-2,-1,0,1,2}
    \draw (\x,-0.07) -- (\x,0.07);
  \node[below, font=\small] at (-4,-0.12) {$-4$};
  \node[below, font=\small] at (0,-0.12) {$0$};
  \draw[omDef, very thick] (-4,0.02) -- (2.5,0.02);
  \draw[omDef, fill=white, thick] (-4,0.02) circle (2.4pt);
\end{tikzpicture}

{\small Penyelesaian $7 - 2x < 15$: semua $x > -4$ (titik berongga:
$-4$ sendiri bukan penyelesaiannya).}
\end{omfigure}

\section{Latihan}

\begin{exercise}[$\star$]\label{exo:g9:algebra:1}
Jabarkan lalu sederhanakan:
\[
4(3x - 2), \qquad
(x + 6)(x + 2), \qquad
(2x - 3)(x + 5), \qquad
(x - 1)(x - 9).
\]
\end{exercise}

\begin{exercise}[$\star$]\label{exo:g9:algebra:2}
Jabarkan dengan memakai kesamaannya:
\[
(x + 8)^2, \qquad
(3x - 4)^2, \qquad
(x + 10)(x - 10), \qquad
(5x + 2)(5x - 2).
\]
\end{exercise}

\begin{exercise}[$\star$]\label{exo:g9:algebra:3}
Hitunglah di kepala, dengan memakai sebuah kesamaan: $101^2$; \quad
$59^2$ (tulislah $59 = 60 - 1$); \quad $53 \times 47$.
\end{exercise}

\begin{exercise}[$\star$]\label{exo:g9:algebra:4}
Faktorkan:
\[
9x + 12, \qquad
x^2 - 64, \qquad
x^2 + 14x + 49, \qquad
18x^2 - 6x .
\]
\end{exercise}

\begin{exercise}[$\star$]\label{exo:g9:algebra:5}
Selesaikan, dengan menulis tiap langkahnya:
\[
4x + 7 = 19, \qquad
6x - 5 = 2x + 11, \qquad
5(x - 3) = 3x + 1 .
\]
\end{exercise}

\begin{exercise}[$\star\star$]\label{exo:g9:algebra:6}
Selesaikan dengan memakai aturan \omterm{def:g2:mult:def}{hasil kali} \omterm{def:g1:counting:zero}{nol}:
\[
(x + 4)(3x - 12) = 0, \qquad
x^2 - 121 = 0, \qquad
(2x - 1)^2 = 0, \qquad
x^2 = 8x .
\]
\end{exercise}

\begin{exercise}[$\star\star$]\label{exo:g9:algebra:7}
Faktorkan $E = (3x + 2)^2 - 25$, lalu selesaikan $E = 0$.
\end{exercise}

\begin{exercise}[$\star\star$]\label{exo:g9:algebra:8}
Selesaikan pertidaksamaannya lalu gambarkan penyelesaiannya pada garis bilangan:
\[
3x - 4 \geq 8, \qquad
5 - 2x > 1, \qquad
4(x + 1) \leq x - 5 .
\]
\end{exercise}

\begin{exercise}[$\star\star$]\label{exo:g9:algebra:9}
Aku memikirkan sebuah bilangan, mengalikannya dengan $3$, menambahkan
$7$, dan memperoleh hasil yang sama seperti kalau aku mengalikannya
dengan $5$ lalu mengurangkan $9$. Berapa bilanganku? (Susunlah sebuah
\omterm{def:g8:equations:def}{persamaan} lalu selesaikan.)
\end{exercise}

\begin{exercise}[$\star\star\star$]\label{exo:g9:algebra:10}
Misalkan $n$ sebuah bilangan bulat.
\begin{enumerate}
  \item Jabarkan $(n + 1)(n - 1)$ lalu simpulkan bahwa \omterm{def:g2:mult:def}{hasil kali}
        kedua tetangga $n$ selalu $n^2 - 1$ (misalnya
        $24 \times 26 = 25^2 - 1 = 624$).
  \item Tunjukkan bahwa \omterm{def:g1:addition:def}{jumlah} tiga bilangan bulat yang berurutan
        selalu \omterm{def:g9:arith:divisor}{kelipatan} $3$.
\end{enumerate}
\end{exercise}

\section{Soal: Aljabar angka}

\begin{problem}[{Soal akhir pekan --- mengapa aturan keterbagian
berhasil, trik 1089 yang dibongkar, dan berhitung di kepala yang
tampak seperti sihir}]
\label{pb:g9:algebra:1}
Sekolah dasar mengajarimu \emph{aturan} (sebuah bilangan \omterm{def:g6:wholes:divisible}{habis dibagi}
$9$ ketika \omterm{def:g1:addition:def}{jumlah} angkanya juga demikian); pesulap memainkan
\emph{trik} (semua orang berakhir di $1089$); pasar mengetahui
\emph{jalan pintas} ($35^2 = 1225$, seketika). Semuanya rahasia yang
sama: bilangan dengan angka $a$, $b$, $c$ adalah bentuk
$100a + 10b + c$, dan aljabar bab ini
(\cref{thm:g9:algebra:identities}, \omterm{def:g9:algebra:expand}{penjabaran}, pemfaktoran) dapat
membongkarnya. Pada akhir soal ini kamu akan sudah membuktikan
aturannya, membongkar triknya, dan mempelajari sihirnya.

\medskip
\noindent\textbf{Bagian I --- Angka di bawah mikroskop.}
\begin{enumerate}
  \item Bilangan dua angka dengan angka puluhan $a$ dan angka satuan
        $b$ adalah $10a + b$. Jabarkan lalu faktorkan \omterm{ex:g1:subtraction:difference}{selisih}
        antara bilangannya dan bilangan baliknya,
        $(10a + b) - (10b + a)$, lalu nyatakan penemuanmu.
        Periksalah pada $72 - 27$.
  \item Buktikan aturan sembilan untuk bilangan tiga angka:
        periksalah kesamaan
        \[
        100a + 10b + c = 9\,(11a + b) + (a + b + c),
        \]
        lalu jelaskan mengapa kesamaan itu menunjukkan bahwa sebuah
        bilangan dan \omterm{def:g1:addition:def}{jumlah} angkanya meninggalkan \emph{\omterm{def:g3:division:remainder}{sisa} yang
        sama} pada pembagian dengan $9$
        --- jadi yang satu \omterm{def:g6:wholes:divisible}{habis dibagi} $9$ persis ketika yang lain
        juga (\cref{prop:g9:arith:rules}, yang kini terbukti).
  \item Simpulkan aturan tiga dari kesamaan yang sama, dalam satu
        kalimat.
  \item ``Membuang sembilan'', pemeriksaan milik para akuntan: untuk
        menguji $47 \times 86 = 4\,042$, gantilah tiap bilangannya
        dengan \omterm{def:g1:addition:def}{jumlah} angka dari \omterm{def:g1:addition:def}{jumlah} angkanya: $47 \to 11 \to 2$,
        $86 \to 14 \to 5$, dan $2 \times 5 = 10 \to 1$; lalu
        $4\,042 \to 10 \to 1$: cocok. Jelaskan mengapa \omterm{def:g2:mult:def}{hasil kali}
        yang benar pasti lulus ujian itu (tulislah $47 = 9k + 2$,
        $86 = 9l + 5$ lalu jabarkan), lalu carilah hasil yang salah
        yang gagal ditangkap ujiannya (mengapa angka yang tertukar
        tidak terlihat olehnya?).
  \item Buktikan aturan sebelas untuk bilangan empat angka:
        periksalah
        \[
        1000a + 100b + 10c + d
        = 11\,(91a + 9b + c) + (- a + b - c + d),
        \]
        lalu pakailah untuk memutuskan apakah $2\,926$ \omterm{def:g6:wholes:divisible}{habis dibagi}
        $11$.
\end{enumerate}

\medskip
\noindent\textbf{Bagian II --- Trik 1089.}
Triknya: pikirkan bilangan tiga angka yang angka pertama dan
terakhirnya berbeda sedikitnya $2$; baliklah urutannya; kurangkan
yang lebih kecil dari yang lebih besar; baliklah urutan hasilnya;
jumlahkan kedua bilangan terakhirnya. Pesulapnya mengumumkan:
\emph{1089}.
\begin{enumerate}[resume]
  \item Mainkan triknya pada $742$ dan pada satu bilangan pilihanmu
        sendiri.
  \item Misalkan bilangannya punya angka $a$, $b$, $c$ dengan
        $a > c$. Faktorkan pengurangan pertamanya,
        $(100a + 10b + c) - (100c + 10b + a)$, lalu sebutkan semua
        nilai yang dapat diambilnya ketika $a - c$ berjalan dari $2$
        sampai $9$.
  \item Masing-masing \omterm{def:g9:arith:divisor}{kelipatan} $99$ itu punya pola angka yang
        mencolok. Dengan menulis $m = a - c$, periksalah kesamaan
        \[
        99m = 100\,(m - 1) + 90 + (10 - m),
        \]
        lalu bacalah ketiga angka $99m$.
  \item Jumlahkan $99m$ dengan bilangan baliknya sendiri, dengan
        memakai angka pada pertanyaan 8, lalu buktikan bahwa
        \omterm{def:g1:addition:def}{jumlahnya} $1089$ --- berapa pun $m$. Triknya mati; hiduplah
        aljabarnya.
  \item Mengapa pesulapnya menuntut angka pertama dan terakhirnya
        berbeda sedikitnya $2$? Periksalah keadaan $a = c$ dan
        $a - c = 1$ (yang pengurangannya memberi $99$, yang harus
        ditulis $099$ agar pembalikannya berhasil), lalu nyatakan
        catatan kecil yang menjaga triknya tetap jujur.
\end{enumerate}

\medskip
\noindent\textbf{Bagian III --- Sihir untuk pemakaian sehari-hari.}
\begin{enumerate}[resume]
  \item Trik ulang tahun: ``kalikan bulan kelahiranmu dengan $5$,
        tambahkan $7$, kalikan dua, tambahkan tanggal kelahiranmu,
        kurangkan $14$.'' Tunjukkan bahwa hasilnya selalu $10m + d$
        --- yaitu bulan dan tanggalnya, terbaca sekilas. (Cobalah
        pada seseorang.)
  \item Kamu merancang ulang triknya menjadi: ``kalikan bulannya
        dengan $5$, tambahkan $9$, kalikan dua, tambahkan
        tanggalnya.'' Bilangan berapa yang kini harus dikurangkan
        penontonnya pada akhirnya agar bacaan bulan--tanggalnya
        sama?
  \item Trik angka yang hilang: penontonnya memilih bilangan apa
        pun, mengurangkan \omterm{def:g1:addition:def}{jumlah} angkanya (pertanyaan 2 berkata
        hasilnya selalu \omterm{def:g9:arith:divisor}{kelipatan} $9$ --- jelaskan mengapa untuk
        bilangan tiga angka), lalu mencoret satu angka yang
        \emph{bukan \omterm{def:g1:counting:zero}{nol}} lalu membacakan angka yang tersisa
        kepadamu. Jelaskan bagaimana kamu memulihkan angka yang
        dicoretnya, dan mengapa $0$ harus dikecualikan.
  \item Sihir pasar: buktikan kesamaan
        $(10t + 5)^2 = 100\,t(t+1) + 25$, nyatakan aturan untuk
        mengkuadratkan bilangan mana pun yang berakhir dengan $5$,
        lalu hitunglah $45^2$, $85^2$ dan $105^2$ di kepalamu.
  \item Penutupnya, yang memperluas \cref{exo:g9:algebra:10}:
        buktikan bahwa \omterm{def:g1:addition:def}{jumlah} lima bilangan bulat berurutan selalu
        \omterm{def:g9:arith:divisor}{kelipatan} $5$, tetapi \omterm{def:g1:addition:def}{jumlah} \emph{empat} bilangan bulat
        berurutan \emph{tidak pernah} \omterm{def:g9:arith:divisor}{kelipatan} $4$.
        (Dua perhitungan pendek --- dan satu pesan moral tentang apa
        yang dilakukan aljabar pada kata ``selalu'' dan ``tidak
        pernah''.)

\end{enumerate}
\end{problem}
